\chapter{Conclusiones}
\label{cha:conclusiones}

Este capítulo contrasta el objetivo general y los objetivos específicos planteados en la Introducción con los resultados medidos sobre el prototipo, e indica para cada uno si se cumplió, si se cumplió en parte o si no se cumplió. A continuación, resume los aportes principales del trabajo, explicita sus limitaciones, incluido lo que el prototipo no llegó a implementar, y presenta las líneas de trabajo futuro que se desprenden de ellas.

\section{Cumplimiento de los objetivos}
\label{sec:conclusiones-objetivos}

El cumplimiento se evalúa en dos pasos: primero cada objetivo específico, contra la cota numérica fijada en la Introducción, y luego el objetivo general, a partir de lo que esos resultados permiten afirmar y de lo que el alcance del prototipo deja sin demostrar.

\subsection{Objetivos específicos}
\label{sec:conclusiones-especificos}

Cada objetivo específico de la Sección~\ref{sec:objetivos-especificos} incluye una cota numérica que el prototipo debía demostrar. La Tabla~\ref{tab:conclusiones-objetivos} compara cada cota con el valor obtenido.

\begin{table}[!htbp]
\centering
\caption{Cumplimiento de los objetivos específicos}
\label{tab:conclusiones-objetivos}
\renewcommand{\arraystretch}{1.3}
\footnotesize
\setlength{\tabcolsep}{4pt}
\begin{tabular}{|>{\centering\arraybackslash}m{0.6cm}|>{\raggedright\arraybackslash}m{4.4cm}|>{\raggedright\arraybackslash}m{5.4cm}|>{\raggedright\arraybackslash}m{2.6cm}|}
\hline
\textbf{N.º} & \textbf{Cota} & \textbf{Resultado} & \textbf{Estado} \\
\hline
% VERIFICAR-COD: los tiempos (fila 1) y las 178 recomendaciones (fila 3) se midieron con seis reglas; volver a medir con las trece
1 & Análisis completo de una tienda en menos de 10 segundos sin el modelo de lenguaje; latencia del panel menor a 300 ms en el 95\% de las consultas & Entre 7,4 y 8,4 segundos (mediana), con un máximo de 9,2 segundos; percentil 95 de la latencia entre 204 y 236 ms & Cumplido \\
\hline
2 & WAPE inferior al 30\% sobre las ventas totales de cada tienda; silueta de al menos 0,4 & WAPE de 16,6\%, 28,1\% y 27,7\%; silueta de 0,662, 0,546 y 0,444 & Cumplido sobre el conjunto generado \\
\hline
3 & 100\% de las recomendaciones con texto explicativo, aun sin el modelo de lenguaje & 178 de 178 recomendaciones con texto: 67 redactadas por el modelo y 111 por la plantilla de respaldo & Cumplido \\
\hline
4 & Cobertura de al menos el 85\% de las líneas del backend y del frontend; pruebas de punta a punta con las dos plataformas & Backend: 85,2\% de las líneas, con 469 pruebas. Frontend: al menos el 85\% de las líneas, con 73 pruebas; pruebas de punta a punta con la tienda de desarrollo de Tiendanube y con una tienda WooCommerce & Cumplido \\
\hline
\end{tabular}

\smallskip
\noindent\textit{Fuente: elaboración propia sobre la base de los resultados de la Sección~\ref{sec:evaluacion-modelos} y de la Sección~\ref{sec:resultados-cobertura}.}
\end{table}


El primer objetivo se cumple con margen en las dos cotas. El análisis completo de una tienda, que comprende el pronóstico, la segmentación y el motor de reglas, tardó entre 7,4 y 8,4 segundos sin la redacción del modelo de lenguaje y nunca superó los 9,2 segundos, y el 95\% de las consultas con las que el panel carga las recomendaciones se respondió en menos de 240 milisegundos (Tablas~\ref{tab:eval-pipeline} y~\ref{tab:eval-latencia}). Con la redacción de Gemini, el análisis tarda entre 41 y 63 segundos; ese tiempo no lo percibe el comerciante, porque el análisis se ejecuta en segundo plano después de cada sincronización. Ambas mediciones corresponden a la escala del conjunto generado y a un backend ejecutado en una computadora personal, por lo que pueden diferir de las de un despliegue con más tiendas.

El segundo objetivo se cumple en los términos en que se formuló, pero el resultado debe leerse junto con la comparación con las referencias. El WAPE sobre las ventas totales quedó por debajo del 30\% en las tres tiendas, y la silueta, por encima de 0,4 (Tablas~\ref{tab:eval-wape} y~\ref{tab:eval-silueta}). Sin embargo, solo en la tienda de gastronomía, de demanda alta y diaria, el modelo supera a las dos referencias ingenuas; en las de indumentaria y electrónica, de demanda intermitente, queda a la par de ellas, y por variante el error supera el 100\%. A su vez, la segmentación recupera los tres perfiles de lealtad que el generador asignó a los clientes. Ambas métricas se midieron sobre datos generados por script, por lo que demuestran que los métodos funcionan sobre series y clientes con la estructura que el generador introdujo, y no que alcancen la misma precisión en una tienda real.

El tercer objetivo se cumple: las 178 recomendaciones generadas en la medición de rendimiento contaron con su texto explicativo. La medición se realizó con la capa gratuita de la API de Gemini, que rechazó por exceso de cuota buena parte de las solicitudes, de modo que el resultado demuestra en condiciones reales el mecanismo de respaldo que el objetivo exigía: cuando el modelo no responde, la recomendación conserva su texto, generado por la plantilla. La priorización según los resultados anteriores, que el mismo objetivo menciona, funciona como se diseñó, aunque solo pudo verificarse sobre un historial simulado (Sección~\ref{sec:eval-bandit}).

El cuarto objetivo se cumple. En el backend, las 469 pruebas automatizadas cubren el 85,2\% de las líneas, por encima de la meta; en el frontend, las 73 pruebas cubren al menos el 85\% de las líneas (Tabla~\ref{tab:resultados-cobertura}). Las pruebas de punta a punta se completaron con las dos plataformas: con la tienda de desarrollo de Tiendanube y con una tienda WooCommerce se probaron la sincronización inicial, la generación de recomendaciones, los \emph{webhooks} y la aplicación dentro del panel de administración.

\subsection{Objetivo general}
\label{sec:conclusiones-general}

El objetivo general planteaba desarrollar un prototipo funcional que transforme los datos operativos de tiendas de comercio electrónico en recomendaciones accionables para emprendedores, con el fin de mejorar su toma de decisiones comerciales. El prototipo cumple la primera parte de ese objetivo: se conecta a tiendas de Tiendanube y de WooCommerce, sincroniza sus productos, variantes, clientes y órdenes, con pruebas de punta a punta en las dos plataformas, y produce recomendaciones concretas, adaptadas al rubro, priorizadas y explicadas en lenguaje natural, que el comerciante consulta dentro del mismo panel de administración en el que opera, % PENDIENTE-IMPL: B1
que en Tiendanube llegan con la acción lista para aplicarse con un clic y su confirmación, y cuyo resultado el sistema evalúa para ajustar la priorización siguiente. Los cuatro objetivos específicos que desagregan esa capacidad se cumplen. El prototipo se encuentra desplegado en Railway, con la API, el proceso en segundo plano y el frontend como servicios separados y la base de datos en una rama de producción de Neon (Sección~\ref{sec:infraestructura}).

La segunda parte, que las recomendaciones mejoren la toma de decisiones de los emprendedores, no puede afirmarse con la evidencia disponible. Demostrarla requiere que comerciantes reales usen el prototipo durante varias semanas y que se observe qué recomendaciones ejecutan y con qué resultado, y ninguna tienda real operó con el prototipo durante el proyecto. Por ese motivo, el objetivo general se considera parcialmente cumplido: el trabajo demuestra que es técnicamente viable producir recomendaciones accionables a partir de los datos de una tienda, pero no todavía que esas recomendaciones mejoren las decisiones de quienes las reciben.

\section{Aportes principales}
\label{sec:conclusiones-aportes}

El primer aporte es la arquitectura del motor de recomendaciones como un sistema híbrido en tres capas (Sección~\ref{sec:motor-ia}): modelos predictivos de pronóstico de demanda y de segmentación de clientes, reglas con umbrales diferenciados por rubro y un bandit con muestreo de Thompson que prioriza los tipos de recomendación según sus resultados. El modelo de lenguaje se limita a redactar el texto de cada recomendación a partir de los números que calcularon las capas anteriores, sin intervenir en qué se recomienda ni en qué orden. Esta separación hace que cada recomendación sea explicable y reproducible a partir de los datos que la originaron, y que una falla del proveedor del modelo de lenguaje degrade solo la redacción y no el análisis.

El segundo aporte es el ciclo de retroalimentación cerrado, ausente en las alternativas relevadas en el Estado del Arte. El sistema registra la ejecución de cada recomendación, ya sea porque el comerciante la aplica desde el panel, la confirma o la descarta indicando que ya la hizo, o porque el sistema la infiere a partir de los cambios de stock y de precio de la tienda, % PENDIENTE-IMPL: B1 (aplicada desde el panel), A4-bandit (descarte «Ya lo hice por mi cuenta»)
 compara el resultado esperado con el real una vez transcurrida la ventana de evaluación y usa ese resultado para actualizar la priorización. La comparación se presenta al comerciante como una correlación indicativa y no como una prueba de causalidad, con lo que el sistema explicita el alcance de lo que mide.

El tercer aporte es la arquitectura multiplataforma (Sección~\ref{sec:multiplataforma}). Un modelo de datos canónico y una interfaz de adaptadores permitieron integrar WooCommerce como segunda plataforma sin modificar el núcleo de análisis, con la misma aplicación dentro del panel de cada plataforma. Esta arquitectura responde a la dependencia de un único proveedor de datos que el análisis estratégico identificó como la principal fuente de riesgo del proyecto, tanto en la matriz FODA cruzado como en el poder de negociación de los proveedores del modelo de Porter, y reduce la incorporación de nuevas plataformas al desarrollo de su adaptador.

El cuarto aporte es el tratamiento de la seguridad y la privacidad como parte del diseño (Sección~\ref{sec:seguridad}): el cifrado con AES-256-GCM de los tokens y credenciales almacenados, el aislamiento estricto entre tiendas, la verificación de la firma de los \emph{webhooks}, la implementación de los derechos de supresión y de acceso previstos en la Ley~25.326 mediante los \emph{webhooks} de privacidad, la anonimización de los datos que se envían al proveedor del modelo de lenguaje y las mitigaciones frente a la inyección de instrucciones a través de los textos de la tienda.

A estos aportes técnicos se suma la evaluación de los modelos con una declaración explícita de sus límites (Sección~\ref{sec:evaluacion-modelos}): el registro del origen de los datos de cada tienda permite informar en todo momento qué proporción de los datos fue generada por script y calcular las métricas solo sobre tiendas reales cuando existan. En el plano del negocio, el análisis de viabilidad indica que el servicio es sostenible en los escenarios base y optimista a partir de 55 suscriptores pagos (Sección~\ref{sec:viabilidad-conclusion}), con la adopción y el precio como variables críticas.

\section{Limitaciones}
\label{sec:conclusiones-limitaciones}

Las conclusiones anteriores están sujetas a las siguientes limitaciones:

\begin{itemize}
    \item \textbf{Datos generados por script.} El 100\% de las órdenes con las que se evaluó el prototipo fue generado por script. Muchos parámetros del generador son supuestos de diseño, fijados para que aparezcan los escenarios de las historias de usuario, y las métricas se miden sobre datos que provienen de los mismos supuestos que los modelos aprenden (Sección~\ref{sec:eval-limitaciones}). La tienda de desarrollo de Tiendanube y la tienda WooCommerce de prueba validan la integración técnica, pero no el comportamiento del motor con ventas reales.
    \item \textbf{Bandit sin uso real.} El bandit aprendió de un historial simulado cuyas tasas de éxito fijó el equipo, por lo que demuestra el mecanismo de actualización y de priorización, pero no qué recomendaciones resultan exitosas en tiendas reales. Las distribuciones previas por rubro incluyen ese historial, y la ejecución inferida del historial se asignó al azar, por lo que tampoco pudo evaluarse la precisión de la inferencia.
    \item \textbf{Muestra del user research.} Las cinco entrevistas surgieron de un muestreo por conveniencia, no incluyen el rubro electrónica y una de ellas, la de E01, está atravesada por un vínculo personal con uno de los autores (Sección~\ref{sec:metodologia-entrevistas}). Sus hallazgos, incluidos los indicios de disposición a pagar, son cualitativos y no generalizables. La sesión de brainstorming validó el diseño desde perfiles técnicos, de producto y de marketing, y no desde una muestra de emprendedores del segmento.
    \item \textbf{Escala y entorno de medición.} Los tiempos se midieron con tres tiendas de entre 17 y 32 variantes con ventas, desde una computadora personal. La mayor parte del tiempo del motor corresponde a escrituras individuales en la base de datos, y la redacción con Gemini se midió en la capa gratuita, que limitó la cantidad de textos redactados por el modelo.
    \item \textbf{Funcionalidades no implementadas.} La vista independiente del ranking de productos más vendidos (HU-04) no se incorporó a la interfaz. La confirmación de la utilidad de una recomendación se infiere, sin una consulta directa al comerciante. La gestión de secretos se realiza mediante variables de entorno, sin rotación automatizada. En WooCommerce, cada tienda se aloja en el servidor de su comerciante, que el backend debe poder alcanzar.
    % PENDIENTE-IMPL: B6 (limitaciones de la acción con un clic)
    \item \textbf{Alcance de la acción con un clic.} En WooCommerce, el prototipo se conecta con claves de solo lectura, por lo que la acción no se aplica con un clic: la recomendación ofrece \enquote{Prefiero hacerlo yo}, con los datos para copiarlos y cargarlos a mano, y la ejecución se registra de forma manual o inferida. La regla de carritos abandonados, a su vez, depende de que la plataforma informe esos carritos, lo que la API de WooCommerce no hace de forma nativa. Por último, cuando una variante no tiene costo cargado, el descuento sugerido no puede respetar el piso de margen y se calcula solo con la tabla por rubro (Tabla~\ref{tab:motor-descuentos}); los porcentajes de esa tabla son, además, supuestos de diseño fundados en las entrevistas, no estimados a partir de resultados.
    \item \textbf{Fecha del último cambio de precio.} La regla de revisión de precio cuenta los días desde el campo \texttt{updated\_at} de la variante, que el sistema registra en cada sincronización en la que la variante llega modificada. Como la plataforma informa cualquier modificación del producto y no solo la del precio, un cambio de stock o de descripción reinicia la cuenta, y al conectar una tienda todas sus variantes empiezan con cero días. Las fotos de \texttt{VARIANT\_SNAPSHOTS}, que ya registran cada cambio de precio, permitirían fecharlo con exactitud.
    \item \textbf{Supuestos del análisis económico.} Los parámetros de adopción y de bajas de los escenarios surgen de indicios cualitativos y no de una medición, y los flujos se calculan antes de impuestos (Sección~\ref{sec:viabilidad-conclusion}).
    \item \textbf{Límites de pedidos detrás del proxy.} La API limita la cantidad de pedidos por minuto que acepta de cada dirección IP. En el despliegue, las solicitudes llegan a la API a través de nginx, y la configuración de proxies de confianza no garantiza que la API distinga la IP original de cada usuario, por lo que usuarios distintos pueden compartir una misma IP y, con ella, el mismo límite de pedidos.
\end{itemize}

\section{Trabajo futuro}
\label{sec:conclusiones-futuro}

Las limitaciones anteriores definen las líneas de trabajo futuro, ordenadas según su relevancia para validar la propuesta:

\begin{enumerate}
    \item \textbf{Validación con tiendas reales.} Conectar la tienda real de un emprendedor del segmento, que opere con el prototipo durante varias semanas con su consentimiento, y aplicar sobre ella el protocolo de la Sección~\ref{sec:eval-limitaciones}: medir el error del pronóstico frente a las referencias y la silueta de la segmentación, medir qué proporción de las recomendaciones se ejecuta y con qué impacto, contrastar la inferencia de la ejecución con la confirmación del comerciante y recalcular las distribuciones previas del bandit sin el historial simulado. Con ese emprendedor, y con los que se sumen después, validar el precio y la conversión del período de prueba que el análisis económico deja como supuestos.
    \item \textbf{User research ampliado.} Repetir las entrevistas con una muestra mayor, convocada de forma abierta y que incluya el rubro electrónica, y realizar una sesión de validación con emprendedores del segmento.
    \item \textbf{Nuevas plataformas y modos de uso.} Integrar Shopify mediante un nuevo adaptador, ofrecer un modo de uso independiente de las plataformas y publicar la aplicación en la tienda de aplicaciones de Tiendanube y en el Woo Marketplace.
    \item \textbf{Infraestructura.} Migrar la infraestructura a AWS, agrupar las escrituras del motor de reglas para reducir su tiempo de ejecución y gestionar los secretos con un servicio dedicado con rotación automatizada.
    \item \textbf{Nuevas reglas y acciones.} Incorporar reglas adicionales: el umbral de envío gratis, que sugiere el monto mínimo de compra a partir del ticket de la tienda; los combos de productos que se compran juntos; la variante ganadora en indumentaria, que concentra la reposición en el color o el talle que más rota; el combo de menú en gastronomía; y la obsolescencia de modelo en electrónica, ante el lanzamiento de una versión nueva. Sumar promociones del tipo 2x1 y 3x2 a las acciones que se aplican con un clic, y aprender también el nivel del descuento con el bandit, con un brazo por tipo de recomendación y nivel de descuento (por ejemplo, 15, 20 y 25\%), en lugar de fijarlo con la tabla por rubro.
    \item \textbf{Funcionalidades.} Incorporar la vista del ranking de productos más vendidos (HU-04); y desarrollar las ideas del brainstorming clasificadas como futuras (Sección~\ref{sec:brainstorming-resultados}): la consulta directa sobre la utilidad de una recomendación, el impacto acumulado de las recomendaciones ejecutadas, un reporte mensual con un resumen semanal de recomendaciones y el seguimiento de acciones realizadas fuera de la plataforma, como las publicaciones en redes sociales.
\end{enumerate}
